\documentclass[10pt,a4paper]{amsart}
\usepackage[margin=19mm]{geometry}
\usepackage{amsmath,amssymb,mathtools}
\usepackage[scr=rsfs,cal=euler]{mathalfa}
\usepackage{xfrac}
\usepackage[unicode,hidelinks,bookmarks=false]{hyperref}
\newcommand{\ie}{i.e.\ }
\newcommand{\cf}{cf.\ }
\newcommand{\resp}{respectively\ }
\newcommand{\Subsection}{\subsection}
\providecommand{\nobreakdash}{\nobreak\hbox{-}}
\setlength{\emergencystretch}{2em}
\multlinegap0pt
\usepackage{amssymb}
\usepackage{cancel}
\usepackage[normalem]{ulem}
\usepackage{tikz}
\newtheorem{theorem}{Theorem}
\newtheorem{lemma}{Lemma}
\newcommand{\Dil}{\operatorname{Dil}}
\DeclareMathOperator{\En}{E}
\newcommand{\Fc}{\mathcal{F}}
\DeclareMathOperator{\HE}{HE}
\newcommand{\Rb}{\mathbb{R}}
\DeclareMathOperator{\TE}{TE}
\newcommand{\Tc}{\mathcal{T}}
\DeclareMathOperator{\tr}{Trace}
\newcommand{\vol}{\operatorname{vol}}
\title{Extremal mappings of tori, Teichm\"uller potentials and symmetric-space distance}
\author{Benson Farb, Eduard Looijenga}
\date{Source: arXiv:2608.15794v1 (CC BY 4.0)}
\begin{document}
\maketitle

\noindent\emph{Source and licence.} Extremal mappings of tori, Teichm\"uller potentials and symmetric-space distance. Version arXiv:2608.15794v1 (CC BY 4.0), distributed by arXiv under the Creative Commons Attribution 4.0 International licence, http://creativecommons.org/licenses/by/4.0/, as recorded in the arXiv licence metadata for this exact version and shown on the abstract page https://arxiv.org/abs/2608.15794v1. Full licence notice: \texttt{licenses/arxiv-2608.15794-CC-BY-4.0.txt}.\par\smallskip

\noindent\textit{Editorial scope.} Counted unit: the article's theorem theorem:newmain (source0.tex lines 313--333). The article's complete original proof of it is delivered with it (source0.tex lines 659--693, one \texttt{proof} environment of 35 lines, which the article places away from the statement). No mathematics is written, repaired or re-derived: the statement and the proof are the article's own lines, in the article's own notation, and no retained line is substituted or re-flowed.  Retained uncounted as the source-local context the proof needs: lines 524--533; lines 592--594; lines 634--639.  Nothing was omitted from any retained span, so the package carries no omission record.  No retained line contains a citation, so the wrapper carries no bibliography and no bibliography entry.  No retained line contains a figure environment, an image-inclusion command or drawing code, so no image is delivered and the package declares no asset dependency.  Editorial formatting only, in addition: an AMS \texttt{amsart} preamble that restates the article's own declarations and macros used by the retained lines, namely the environments \texttt{theorem}, \texttt{lemma} on separate counters, together with the standard packages those lines need.  The retained environments print in this document as Theorem 1, Lemma 1 in order of appearance, whereas the article prints the same items with the numbering of its own full document.  The complete native source of the article as distributed in the arXiv e-print for this exact version is bundled unchanged as \texttt{sources/207784/source0.tex}.
\begin{theorem}[{\bf Affine maps minimize $\Dil_k$ and $\En_k$ in their homotopy class}]
\label{thm:extremal}
Let $\Tc_0$ and $\Tc_1$ be flat  tori, not necessarily of the same dimension.  Let $\varphi: H_1(\Tc_0)\to H_1(\Tc_1)$ be a linear map and denote 
by $\Fc_\varphi$ the associated homotopy class of maps $\Tc_0\to \Tc_1$.
Then for any Lipschitz map 
$f\in \Fc_\varphi$, any affine $\Phi\in \Fc_\varphi$ and any $k\ge 1$, 
\[\Dil_k(f)\geq \Dil_k(\Phi)=\|\wedge^k\varphi\| \quad \text{and}\quad \En_k(f)\ge  \En_k(\Phi)=\|\wedge^k\varphi\|_{\HE}^2\vol(\Tc_0).
\]
Furthermore, if $\En_1(f)= \En_1(\Phi)$ then $f$ differs from $\Phi$ by a translation (hence is affine).
\end{theorem}

\begin{equation}\label{eqn:average}
\textstyle \int_{\Tc_0} \wedge^k Df\,d\!\vol_{\Tc_0}=\int_{\Tc_0} \wedge^k D\Phi\,d\!\vol_{\Tc_0}=\vol(\Tc_0).\wedge^k\varphi .
\end{equation}

\begin{lemma}\label{lemma:linalg}
Let $E$ and $F$ be real inner product  spaces of the same dimension $n$ and let $A: E\to F$ and $B: E\to F$ be linear maps with $A$ an isomorphism.  Suppose 
that $\|\wedge^k B\|\le  \|\wedge^kA\|$ for $k=1, \dots, n$. Then 
\[\tr(A^*B) \le \tr(A^*A)\] 
and if equality holds then $B=A$.
\end{lemma}

\begin{theorem}[{\bf Main Theorem}]
\label{theorem:newmain}
Let $\Tc_0$ and $\Tc_1$ be flat tori  of the same dimension $n\ge 1$. 
 Assume  that $d:=\vol(\Tc_0)/\vol(\Tc_1)$ is a positive integer and let $\Fc$ be a homotopy class of maps $\Tc_0\to \Tc_1$ of absolute degree $d$.  
 \medskip
 
\begin{enumerate}
\item The functions  $\Dil_k$, $\TE$, $\En_k$ and $\HE$   have a minimum  on the set  of  \emph{volume-preserving} Lipschitz 
maps  $f$ in $\Fc$  (\footnote{By ``volume-preserving'' we mean $|Jac_f(x)|=1$ 
almost everywhere.  Note that Lipschitz maps $f$ are differentiable almost everywhere by Rademacher's Theorem.}) and take that  minimum on the space of affine maps  in $\Fc$. For $\TE$ and $\HE$ this minimum is taken on affine maps \emph{only} (\footnote{It can be shown that this is also true for  $\En_1, \dots, \En_{n-1}$.}).
\bigskip

\item[(2)] Both Teichm\"uller potentials $\TE$ and $\HE$  induce the symmetric space 
distance on $X_n$:  if  $(\Tc_0,\Fc_0)$ and $(\Tc_1,\Fc_1)$ are marked flat unit volume $n$-tori
then 
\begin{equation}\label{eqn:TEmin}
\textstyle d_{X_n}((\Tc_0,\Fc_0),(\Tc_1,\Fc_1))=\inf_{h\in \Fc_1\circ \Fc_0^{-1}}\TE(h)=\inf_{h\in \Fc_1\circ \Fc_0^{-1}}\HE(h)
\end{equation}
where the infimum is taken over all volume-preserving Lipschitz maps $h:\Tc_0\to \Tc_1$ in the homotopy class $\Fc_1\circ \Fc_0^{-1}$.
\end{enumerate}
\end{theorem}

\begin{proof}[Proof of Part~(1) of Theorem \ref{theorem:newmain}]
Let $f:  \Tc_0\to \Tc_1$  be  a volume-preserving Lipschitz map of absolute degree $d=\vol(\Tc_0)/\vol(\Tc_1)$ and let $\varphi: H_1(\Tc_0)\to H_1(\Tc_1)$ the map induced on homology.  Let 
$a_1\ge a_2\ge\cdots \ge a_n>0$ be the singular sequence of 
$\varphi: H_1(\Tc_0; \Rb)\to H_1(\Tc_1; \Rb)$, and let $\Phi:\Tc_0\to\Tc_1$ be an affine map homotopic to $f$. 

We put $\lambda_k:=\log a_k$. So $\lambda_k\ge \lambda_{k+1}$ and 
$\log \Dil_k(\Phi)=\lambda_1+\cdots +\lambda_k$. Since  $f$ and $\Phi$ are volume-preserving, 
\[\Dil_n(f)=\Dil_n(\Phi)=a_1a_2\cdots a_n=1\] 
and $\Dil_0(f)=\Dil_0(\Phi)=1$ by definition. 
By Theorem~\ref{thm:extremal}, $\Dil_k(f)\ge \Dil_k(\Phi)$ and so $t_k:=\log  \Dil_k(f)-\log \Dil_k(\Phi) \ge 0$.  Then
\begin{multline*}
\textstyle     \TE(f)^2-\TE(\Phi)^2 =\\
\textstyle = \sum_{k=1}^n (\log D_k(f)-\log D_{k-1}(f) )^2 -  \sum_{k=1}^n (\log D_k(\Phi)-\log D_k(\Phi))^2=\\
\textstyle =\sum_{k=1}^n  \bigl( (\lambda_k+t_k-t_{k-1})^2 -\lambda_k^2  \bigr) =\sum_{k=1}^n (t_k-t_{k-1})^2  +
2\sum_{k=1}^n\lambda_k(t_k-t_{k-1}) 
\end{multline*}
and since Abel summation gives  
\[
\textstyle  2\sum_{k=1}^n\lambda_k(t_k-t_{k-1})=2\sum_{k=1}^{n-1}(\lambda_k -\lambda_{k+1})t_k\ge 0,
\]
 this proves that $\TE(f)\geq \TE(\Phi)$. 

 
Suppose now that $\TE(f)=\TE(\Phi)$. Then the preceding shows that $\sum_{k=1}^n(t_k-t_{k-1})^2=0$ and hence $t_k=t_{k-1}$ for every $k$. Since 
$t_0=0$, it follows that $t_k=0$ for all $k$, so that $\Dil_k(f)=\Dil_k(\Phi)$ for $k=1, \dots ,n$.
We regard the derivatives $Df_p$ (where defined) and $d\Phi_p$ as maps between the tangent spaces of $\Tc_0$ and $\Tc_1$ at their identity element (so that now $d\Phi_p=\varphi$).  For almost every $p\in \Tc_0$ where $f$ is smooth, Lemma \ref{lemma:linalg} with $A=\varphi$ and $B=Df_p$ yields
\begin{equation}\label{eqn:trace-pointwise}
\tr (\varphi^*Df_p)\le \tr (\varphi^*\varphi)
\end{equation}
or  equivalently,  $\tr (\varphi^*(\varphi-Df_p))\ge 0$. The identity   \eqref{eqn:average}  gives that  $\int_{\Tc_0} (Df-\varphi)  d\!\vol_{\Tc_0}=0$
and so the integral of the nonnegative function $\tr (\varphi^*(\varphi-Df_p))$ is zero. This implies that $\tr (\varphi^*(\varphi-Df_p))$ is zero almost everywhere. Invoking Lemma \ref{lemma:linalg} once more, we then find that $Df_p$ is constant equal to 
$\varphi$ almost everywhere. This implies that $f$ differs from $\Phi$ by a translation, hence is affine. 

The remaining assertion,  that  $\HE(f)=\HE(\Phi)$ yields the same conclusion, will be proved later.
\end{proof}

\end{document}
